% Printed 441.
\par\clearpage
\noindent\null\hfill{\small\textbf{HÖHENMESSUNGEN.}}\hfill\makebox[0pt][r]{\fontsize{13}{16}\selectfont 441}\par\medskip
\par\noindent\hspace*{\parindent}Das ganze System meiner diesjährigen Messungen umfasst 29 Linien;
zwei davon gehören aber, was die Messung der gegenseitigen Zenithdistanzen
betrifft, schon zu den Arbeiten von 1822 und 1823 (nemlich Falkenberg-Wilsede
und Wilsede-Hamburg). Bei zwei Linien ist die Zenithdistanz \textls[120]{bloss} einseitig
gemessen (nemlich gar nicht von Bremen nach Bottel und von Brüttendorf
nach Zeven); bei zwei andern ist sie auch insofern nur einseitig gemessen,
als bei der Messung von der andern Seite, aus Mangel an Zeit, bloss der
Theodolith angewandt ist (nemlich von Bottel nach Steinberg und von Wilsede
nach Bullerberg\textsuperscript{*}\,). Endlich bei der Linie Litberg-Hamburg\textsuperscript{**} ist in Hamburg
der Knopf der Zielpunkt gewesen, dessen Höhe über den Fenstern des Cabinets ich noch nicht genau kenne. Es bleiben also noch 22 Linien übrig, wo
die reciproken Zenithdistanzen vollständig gemessen sind, d.~i. mit dem 12-zölligen Kreise und unter wenigstens 20-maliger Repetition, und wo \textls[120]{ohne Ausnahme} Heliotroplicht der Zielpunkt gewesen ist. Die Summe der 22 Krümmungen (Amplituden) ist \(6^\circ15'49^{\prime\prime}844\) (die grösste darunter von Wilsede nach
Steinberg \(26'4^{\prime\prime}{,}883\)); die Summe der 22 Refractionen hingegen \(55'32^{\prime\prime}{,}329\).
Also das Mittelverhältniss wie 1 zu
\[0{,}14778.\]
\par\noindent\hspace*{\parindent}Das Mittel aus den einzelnen 22 Quotienten ist \(0{,}14499\); ich ziehe aber
jenes Resultat vor, weil bei \textls[120]{kleinen} Bogen der Quotient viel mehr schwankt
als bei grössern. Ich habe schon anderswo bemerkt, dass ich unter Refraction
die Verschiedenheit der Richtungen an den beiden Endpunkten verstehe; die
meisten Schriftsteller nennen sonst Refraction die Hälfte jener Verschiedenheit,
nemlich die Winkel der Tangenten am Wege des Lichtstrahls an den beiden
Endpunkten mit der Chorde. So verstanden geben also meine diesjährigen
Messungen \(0{,}07389\). Meine sämmtlichen zu diesem Zweck brauchbaren Messungen von 1821---1823 hatten für 34 Linien die Summe der Krümmungen
\(10^\circ8'14^{\prime\prime}{,}249\), die Summe der Refractionen \(=1^\circ17'43^{\prime\prime}{,}670\), also das Verhältniss wie 1 zu \(0{,}12779\) gegeben. Die grösste Linie darunter ist Lichtenberg-Falkenberg, deren Krümmung \(46'7^{\prime\prime}{,}250\) beträgt. Man würde sich sehr irren,
wenn man glaubte, dass der grössere Quotient von 1824 den flächern Gegenden eigen wäre. In der That geben die Messungen von 1823, die ebenso
flachen Gegenden angehören, sogar noch ein kleineres Resultat, als alle von
\par\medskip\noindent\hspace*{0.03\linewidth}{\footnotesize IX.}\hfill\makebox[0.18\linewidth][l]{56}\par

% Printed 442.
\par\clearpage
\noindent\makebox[0pt][l]{\fontsize{13}{16}\selectfont 442}\hfill{\small\textbf{BRIEFWECHSEL.}}\hfill\null\par\medskip
\noindent 1821---1823. Ich bin vielmehr über die Quelle des Unterschiedes gar nicht
zweifelhaft. In den Jahren von 1821---1823 habe ich die Zenithdistanzen fast
alle Vormittags oder Mittags oder bald nach Mittag gemessen, da ich diese Messungen als etwas untergeordnetes betrachtete und in der Regel nur die Stunden
dazu verwandte, wo das Sehen für die Theodolithen-Messungen nicht gut
genug war. Das sind aber an \textls[120]{sonnigen Tagen} ohne Ausnahme die genannten Stunden in flachen Gegenden, und dann ist ebenso beständig an
sonnigen Tagen die Refraction allemal kleiner, als in den etwas spätern Nachmittagsstunden, wo die Bilder ruhiger und schärfer werden. Im Jahr 1824
gingen aber fast die meisten meiner Linien so knapp über zwischenliegende
Hindernisse weg, dass Vormittags bei sehr vielen gar kein Licht herüber
konnte, und ich daher mit wenigen Ausnahmen die Vormittags-Messungen bei
Hauptrichtungen ganz aufgab. Die Zenithdistanzen sind daher 1824 \textls[120]{sämmtlich} Nachmittags bei schon besserer Luft gemessen, doch nie kurz vor Sonnenuntergang, wo nach allen meinen Erfahrungen die Refraction noch bedeutend
grösser wird. Meistens sind sie Nachmittags von 3---4 Uhr gemessen\textsuperscript{*}), und
ich habe so viel als thunlich bei jeder Linie die Messung der gegenseitigen
Zenithdistanzen ungefähr unter gleicher Luftbeschaffenheit zu machen gesucht.
\par\noindent\hspace*{\parindent}Wenn Sie alle meine \(22+34=56\) Resultate einzeln zu haben wünschen,
so stehen sie Ihnen gern zu Dienste. Bei den ausgeschlossenen oben mit \textsuperscript{* **}
bezeichneten Linien sind übrigens
\begin{center}
\setlength{\tabcolsep}{4pt}
\begin{tabular}{lrrr}
& die Krümmungen & die Refractionen & Quotient\\
{}[Wilsede-Bullerberg]\textsuperscript{*} &\(18'56^{\prime\prime}{,}002\)&\(2'42^{\prime\prime}{,}975\)&\(0{,}14346\)\\
{}[Litberg-Hamburg]\textsuperscript{**} &\(16\phantom{' }14{,}047\)&\(2\phantom{' }49{,}295\)&\(0{,}17380\)\\\cline{2-4}
\multicolumn{1}{r}{Summa}&\(35'10^{\prime\prime}{,}049\)&\(5'32^{\prime\prime}{,}270\)&\(0{,}15747\),
\end{tabular}
\end{center}
\noindent indem ich für \textsuperscript{**} die Höhe des Knopfs in Hamburg über den Fenstern
\(=15{,}558\,\mathrm{m}\) annehme; durch eine genauere Bestimmung dieser relativen Höhen
und der Ihres Barometers in Altona gegen obige beiden Punkte werden Sie
mich verbinden.\quad\mbox{.\ .\ .\ .\ .\ .}
\par\medskip\noindent\rule{.25\linewidth}{.4pt}\par
{\small\noindent\hspace*{\parindent}*) So lange die Tage noch länger waren; späterhin successive etwas früher, im October wallte in
der Regel die Luft schon um \(1\frac{1}{2}\) oder 2 Uhr nicht viel mehr.\par}
